\documentclass[12pt]{ximera}
\typeout{Start loading xmPreamble.tex}%

% Add here extra macro's that are loaded automatically by all documents of claas 'ximera' or 'xourse' in this repo

%%
%%  Example:
%%
% \newcommand{\R}{\mathbb{R}}
\usepackage{tikz}
\usetikzlibrary{positioning, arrows.meta}
\usepackage{multicol}
\usepackage{enumitem}
\usepackage{caption}
\usepackage{amsmath}
\usepackage{amsthm}
\usepackage{fontenc}

%% ---------------------------------------------------------------
%% Shared worksheet macros.
%%
%% These came from the Summer 2026 worksheet set, where each worksheet carried its
%% own copy. They have to live here instead: a xourse inlines every activity into a
%% single document, so duplicate \newcommand definitions across activities collide,
%% and the xourse gobbles \input so a per-topic macro file never loads.
%%
%% They use \providecommand, NOT \newcommand. ximera.cls loads this file automatically,
%% and every activity also carries an explicit \typeout{Start loading xmPreamble.tex}%

% Add here extra macro's that are loaded automatically by all documents of claas 'ximera' or 'xourse' in this repo

%%
%%  Example:
%%
% \newcommand{\R}{\mathbb{R}}
\usepackage{tikz}
\usetikzlibrary{positioning, arrows.meta}
\usepackage{multicol}
\usepackage{enumitem}
\usepackage{caption}
\usepackage{amsmath}
\usepackage{amsthm}
\usepackage{fontenc}

%% ---------------------------------------------------------------
%% Shared worksheet macros.
%%
%% These came from the Summer 2026 worksheet set, where each worksheet carried its
%% own copy. They have to live here instead: a xourse inlines every activity into a
%% single document, so duplicate \newcommand definitions across activities collide,
%% and the xourse gobbles \input so a per-topic macro file never loads.
%%
%% They use \providecommand, NOT \newcommand. ximera.cls loads this file automatically,
%% and every activity also carries an explicit \typeout{Start loading xmPreamble.tex}%

% Add here extra macro's that are loaded automatically by all documents of claas 'ximera' or 'xourse' in this repo

%%
%%  Example:
%%
% \newcommand{\R}{\mathbb{R}}
\usepackage{tikz}
\usetikzlibrary{positioning, arrows.meta}
\usepackage{multicol}
\usepackage{enumitem}
\usepackage{caption}
\usepackage{amsmath}
\usepackage{amsthm}
\usepackage{fontenc}

%% ---------------------------------------------------------------
%% Shared worksheet macros.
%%
%% These came from the Summer 2026 worksheet set, where each worksheet carried its
%% own copy. They have to live here instead: a xourse inlines every activity into a
%% single document, so duplicate \newcommand definitions across activities collide,
%% and the xourse gobbles \input so a per-topic macro file never loads.
%%
%% They use \providecommand, NOT \newcommand. ximera.cls loads this file automatically,
%% and every activity also carries an explicit \input{../xmPreamble}, so this file is
%% read twice. That was harmless while it held only \usepackage lines (LaTeX skips a
%% package it has already loaded) but a second \newcommand is a hard error.
%%
%%   \blank[width]        fill-in-the-blank rule (default 2.5cm)
%%   \showwork            "show and explain your work" reminder
%%   \showworkline        ... plus which schedule line was used
%%   \showworkyear        ... plus which year's schedule was used
%%   \schedhead{...}      centred bold caption above a tiered-rate table
%%   \samesquares[scale]{left}{right}   two equal-area squares, labelled
%%   \samecubes[scale]{left}{right}     two equal-volume cubes, labelled
%% ---------------------------------------------------------------

\providecommand{\blank}[1][2.5cm]{\underline{\hspace{#1}}}
\providecommand{\showwork}{\textbf{REMEMBER TO SHOW AND EXPLAIN YOUR WORK.}}
\providecommand{\showworkline}{\textbf{REMEMBER TO SHOW AND EXPLAIN YOUR WORK.
  Explanation should include how and why you know which line of the
  schedule was used to calculate this amount.}}
\providecommand{\showworkyear}{\begin{sloppypar}\textbf{REMEMBER TO SHOW AND
  EXPLAIN YOUR WORK. Explanation should include how and why you know which
  line of the year's tax schedule was used to calculate this tax
  amount.}\end{sloppypar}}
\providecommand{\schedhead}[1]{\begin{center}\textbf{#1}\end{center}}

\providecommand{\samesquares}[3][1]{%
\begin{center}
{\footnotesize These squares are the \textbf{same size}.}

\vspace{1.5mm}
\begin{tikzpicture}[scale=#1,every node/.style={scale=#1}]
  \draw (0,0) rectangle (2.4,2.4);
  \node[anchor=north west] at (0.15,2.25) {Area:};
  \node[anchor=east]  at (-0.15,1.2) {#2};
  \node[anchor=north] at (1.2,-0.15) {#2};
  \begin{scope}[xshift=4.6cm]
    \draw (0,0) rectangle (2.4,2.4);
    \node[anchor=north west] at (0.15,2.25) {Area:};
    \node[anchor=east]  at (-0.15,1.2) {#3};
    \node[anchor=north] at (1.2,-0.15) {#3};
  \end{scope}
\end{tikzpicture}
\end{center}}

\providecommand{\samecubes}[3][1]{%
\begin{center}
{\footnotesize These cubes are the \textbf{same size}.}

\vspace{1.5mm}
\begin{tikzpicture}[scale=#1,every node/.style={scale=#1}]
  \foreach \dx in {0,5.2} {
    \begin{scope}[xshift=\dx cm]
      \draw (0,0) rectangle (2.0,2.0);
      \draw (0,2.0) -- (0.65,2.6) -- (2.65,2.6) -- (2.0,2.0);
      \fill[black!12] (2.0,0) -- (2.65,0.6) -- (2.65,2.6) -- (2.0,2.0) -- cycle;
      \draw (2.0,0) -- (2.65,0.6) -- (2.65,2.6) -- (2.0,2.0);
      \node[anchor=north west] at (0.1,1.9) {Volume:};
    \end{scope}
  }
  \node[anchor=east]  at (-0.15,1.0) {#2};
  \node[anchor=north] at (1.0,-0.15) {#2};
  \node[anchor=west]  at (2.25,0.4) {#2};
  \node[anchor=east]  at (5.05,1.0) {#3};
  \node[anchor=north] at (6.2,-0.15) {#3};
  \node[anchor=west]  at (7.45,0.4) {#3};
\end{tikzpicture}
\end{center}}

%% NOTE: do NOT add \usepackage to individual activity .tex files.
%% When a xourse inlines an activity it gobbles \documentclass and \input, but a bare
%% \usepackage still executes -- after \begin{document} -- and errors out.
%% All shared packages and macros belong here.
%%
%% ximera.cls already loads: amsmath, amssymb, amsthm, cancel, enumitem, graphicx,
%% hyperref, listings, pgfplots, tikz, url, xcolor. No need to re-load those.
, so this file is
%% read twice. That was harmless while it held only \usepackage lines (LaTeX skips a
%% package it has already loaded) but a second \newcommand is a hard error.
%%
%%   \blank[width]        fill-in-the-blank rule (default 2.5cm)
%%   \showwork            "show and explain your work" reminder
%%   \showworkline        ... plus which schedule line was used
%%   \showworkyear        ... plus which year's schedule was used
%%   \schedhead{...}      centred bold caption above a tiered-rate table
%%   \samesquares[scale]{left}{right}   two equal-area squares, labelled
%%   \samecubes[scale]{left}{right}     two equal-volume cubes, labelled
%% ---------------------------------------------------------------

\providecommand{\blank}[1][2.5cm]{\underline{\hspace{#1}}}
\providecommand{\showwork}{\textbf{REMEMBER TO SHOW AND EXPLAIN YOUR WORK.}}
\providecommand{\showworkline}{\textbf{REMEMBER TO SHOW AND EXPLAIN YOUR WORK.
  Explanation should include how and why you know which line of the
  schedule was used to calculate this amount.}}
\providecommand{\showworkyear}{\begin{sloppypar}\textbf{REMEMBER TO SHOW AND
  EXPLAIN YOUR WORK. Explanation should include how and why you know which
  line of the year's tax schedule was used to calculate this tax
  amount.}\end{sloppypar}}
\providecommand{\schedhead}[1]{\begin{center}\textbf{#1}\end{center}}

\providecommand{\samesquares}[3][1]{%
\begin{center}
{\footnotesize These squares are the \textbf{same size}.}

\vspace{1.5mm}
\begin{tikzpicture}[scale=#1,every node/.style={scale=#1}]
  \draw (0,0) rectangle (2.4,2.4);
  \node[anchor=north west] at (0.15,2.25) {Area:};
  \node[anchor=east]  at (-0.15,1.2) {#2};
  \node[anchor=north] at (1.2,-0.15) {#2};
  \begin{scope}[xshift=4.6cm]
    \draw (0,0) rectangle (2.4,2.4);
    \node[anchor=north west] at (0.15,2.25) {Area:};
    \node[anchor=east]  at (-0.15,1.2) {#3};
    \node[anchor=north] at (1.2,-0.15) {#3};
  \end{scope}
\end{tikzpicture}
\end{center}}

\providecommand{\samecubes}[3][1]{%
\begin{center}
{\footnotesize These cubes are the \textbf{same size}.}

\vspace{1.5mm}
\begin{tikzpicture}[scale=#1,every node/.style={scale=#1}]
  \foreach \dx in {0,5.2} {
    \begin{scope}[xshift=\dx cm]
      \draw (0,0) rectangle (2.0,2.0);
      \draw (0,2.0) -- (0.65,2.6) -- (2.65,2.6) -- (2.0,2.0);
      \fill[black!12] (2.0,0) -- (2.65,0.6) -- (2.65,2.6) -- (2.0,2.0) -- cycle;
      \draw (2.0,0) -- (2.65,0.6) -- (2.65,2.6) -- (2.0,2.0);
      \node[anchor=north west] at (0.1,1.9) {Volume:};
    \end{scope}
  }
  \node[anchor=east]  at (-0.15,1.0) {#2};
  \node[anchor=north] at (1.0,-0.15) {#2};
  \node[anchor=west]  at (2.25,0.4) {#2};
  \node[anchor=east]  at (5.05,1.0) {#3};
  \node[anchor=north] at (6.2,-0.15) {#3};
  \node[anchor=west]  at (7.45,0.4) {#3};
\end{tikzpicture}
\end{center}}

%% NOTE: do NOT add \usepackage to individual activity .tex files.
%% When a xourse inlines an activity it gobbles \documentclass and \input, but a bare
%% \usepackage still executes -- after \begin{document} -- and errors out.
%% All shared packages and macros belong here.
%%
%% ximera.cls already loads: amsmath, amssymb, amsthm, cancel, enumitem, graphicx,
%% hyperref, listings, pgfplots, tikz, url, xcolor. No need to re-load those.
, so this file is
%% read twice. That was harmless while it held only \usepackage lines (LaTeX skips a
%% package it has already loaded) but a second \newcommand is a hard error.
%%
%%   \blank[width]        fill-in-the-blank rule (default 2.5cm)
%%   \showwork            "show and explain your work" reminder
%%   \showworkline        ... plus which schedule line was used
%%   \showworkyear        ... plus which year's schedule was used
%%   \schedhead{...}      centred bold caption above a tiered-rate table
%%   \samesquares[scale]{left}{right}   two equal-area squares, labelled
%%   \samecubes[scale]{left}{right}     two equal-volume cubes, labelled
%% ---------------------------------------------------------------

\providecommand{\blank}[1][2.5cm]{\underline{\hspace{#1}}}
\providecommand{\showwork}{\textbf{REMEMBER TO SHOW AND EXPLAIN YOUR WORK.}}
\providecommand{\showworkline}{\textbf{REMEMBER TO SHOW AND EXPLAIN YOUR WORK.
  Explanation should include how and why you know which line of the
  schedule was used to calculate this amount.}}
\providecommand{\showworkyear}{\begin{sloppypar}\textbf{REMEMBER TO SHOW AND
  EXPLAIN YOUR WORK. Explanation should include how and why you know which
  line of the year's tax schedule was used to calculate this tax
  amount.}\end{sloppypar}}
\providecommand{\schedhead}[1]{\begin{center}\textbf{#1}\end{center}}

\providecommand{\samesquares}[3][1]{%
\begin{center}
{\footnotesize These squares are the \textbf{same size}.}

\vspace{1.5mm}
\begin{tikzpicture}[scale=#1,every node/.style={scale=#1}]
  \draw (0,0) rectangle (2.4,2.4);
  \node[anchor=north west] at (0.15,2.25) {Area:};
  \node[anchor=east]  at (-0.15,1.2) {#2};
  \node[anchor=north] at (1.2,-0.15) {#2};
  \begin{scope}[xshift=4.6cm]
    \draw (0,0) rectangle (2.4,2.4);
    \node[anchor=north west] at (0.15,2.25) {Area:};
    \node[anchor=east]  at (-0.15,1.2) {#3};
    \node[anchor=north] at (1.2,-0.15) {#3};
  \end{scope}
\end{tikzpicture}
\end{center}}

\providecommand{\samecubes}[3][1]{%
\begin{center}
{\footnotesize These cubes are the \textbf{same size}.}

\vspace{1.5mm}
\begin{tikzpicture}[scale=#1,every node/.style={scale=#1}]
  \foreach \dx in {0,5.2} {
    \begin{scope}[xshift=\dx cm]
      \draw (0,0) rectangle (2.0,2.0);
      \draw (0,2.0) -- (0.65,2.6) -- (2.65,2.6) -- (2.0,2.0);
      \fill[black!12] (2.0,0) -- (2.65,0.6) -- (2.65,2.6) -- (2.0,2.0) -- cycle;
      \draw (2.0,0) -- (2.65,0.6) -- (2.65,2.6) -- (2.0,2.0);
      \node[anchor=north west] at (0.1,1.9) {Volume:};
    \end{scope}
  }
  \node[anchor=east]  at (-0.15,1.0) {#2};
  \node[anchor=north] at (1.0,-0.15) {#2};
  \node[anchor=west]  at (2.25,0.4) {#2};
  \node[anchor=east]  at (5.05,1.0) {#3};
  \node[anchor=north] at (6.2,-0.15) {#3};
  \node[anchor=west]  at (7.45,0.4) {#3};
\end{tikzpicture}
\end{center}}

%% NOTE: do NOT add \usepackage to individual activity .tex files.
%% When a xourse inlines an activity it gobbles \documentclass and \input, but a bare
%% \usepackage still executes -- after \begin{document} -- and errors out.
%% All shared packages and macros belong here.
%%
%% ximera.cls already loads: amsmath, amssymb, amsthm, cancel, enumitem, graphicx,
%% hyperref, listings, pgfplots, tikz, url, xcolor. No need to re-load those.

\title{Homework 4: Taxes Across Three Eras}
\begin{document}
\begin{abstract}
Comparing the federal income tax burden on inflation-equivalent incomes in 1963, 1988,
and 2013 --- first computing tax forwards from income, then working backwards from a
fixed \$3,000 tax bill, and interpreting what the effective rates reveal once inflation
is accounted for.
\end{abstract}
\maketitle

\begin{instructorNotes}
\textbf{This activity is not yet complete.} Both parts below direct students to ``the
tax schedules provided,'' but the 1963, 1988, and 2013 schedules for married-filing-
jointly are not currently included anywhere in this activity. Until they are added, the
numerical questions cannot be answered. The structure and the reasoning questions are
all here and ready; only the three rate tables are missing.
\end{instructorNotes}

\section*{Part 1 --- More with Taxes}

Throughout, assume a couple who are married filing jointly with no dependents, using
the standard deduction for that year.

\begin{center}
\begin{tabular}{|l|c|c|c|}
\hline
\textbf{Description} & \textbf{1963} & \textbf{1988} & \textbf{2013} \\
\hline
Adjusted Gross Income (AGI) & \$10,500 & \$40,600 & \$80,000 \\
\hline
Standard deduction & \$1,200 & \$8,900 & \$20,000 \\
\hline
Taxable income & & & \\
\hline
Total tax (from schedule) & & & \\
\hline
Effective rate (\% of TOTAL income) & & & \\
\hline
\end{tabular}
\end{center}

\begin{problem}
Use the tax information from \textbf{1963} to calculate how much tax this couple would
pay on \$10,500 of annual income, given a standard deduction of \$1,200. Record the
taxable income and total tax for the 1963 column.

\showwork

\begin{freeResponse}
\end{freeResponse}

\begin{solution}
Taxable income is $\$10{,}500 - \$1{,}200 = \$9{,}300$. Locate that in the 1963
married-filing-jointly schedule, then apply that row's instruction: the base amount plus
the row's rate on the amount over the row's lower boundary.
\end{solution}
\end{problem}

\begin{problem}
Do the same using the \textbf{1988} information: \$40,600 of income, \$8,900 standard
deduction.

\showwork

\begin{freeResponse}
\end{freeResponse}

\begin{solution}
Taxable income is $\$40{,}600 - \$8{,}900 = \$31{,}700$, then apply the 1988 schedule.
\end{solution}
\end{problem}

\begin{problem}
Do the same using the \textbf{2013} information: \$80,000 of income, \$20,000 standard
deduction.

\showwork

\begin{freeResponse}
\end{freeResponse}

\begin{solution}
Taxable income is $\$80{,}000 - \$20{,}000 = \$60{,}000$, then apply the 2013 schedule.
\end{solution}
\end{problem}

\begin{problem}
Calculate the effective tax rate for each year in the table. Calculate it as a
percentage of the \textbf{TOTAL} income (AGI), not of the taxable income.

\begin{freeResponse}
\end{freeResponse}

\begin{solution}
For each year, divide that year's total tax by that year's AGI --- \$10,500, \$40,600,
and \$80,000 respectively --- and express the result as a percentage. Using AGI as the
denominator throughout is what makes the three years comparable, since the standard
deductions differ enormously across them.
\end{solution}
\end{problem}

According to a U.S.\ inflation calculator
(\url{http://www.usinflationcalculator.com/}), all of the adjusted gross incomes above
are roughly equivalent when adjusted for inflation. That is, earning \$80,000 in 2013 is
roughly equivalent to earning \$40,600 in 1988, which is roughly equivalent to earning
\$10,500 in 1963.

\begin{problem}
In which year was this level of income taxed proportionally the most? Explain your
choice.

\begin{freeResponse}
\end{freeResponse}

\begin{solution}
Compare the three effective rates you computed and name the largest. The explanation
should point at the effective rate specifically --- not at the raw dollar amount of tax,
which is not comparable across years, and not at the top marginal rate of each schedule,
which describes a bracket this couple never reaches.
\end{solution}
\end{problem}

\begin{problem}
Why do we need to know that these incomes are all ``roughly equivalent'' when adjusted
for inflation in order to answer the previous part?

\begin{freeResponse}
\end{freeResponse}

\begin{solution}
Because without that fact, the three columns would describe three different standards of
living, and any difference in effective rate could simply be the tax system treating a
richer household differently --- which is what a progressive system is designed to do.

Knowing the incomes are inflation-equivalent means all three columns describe
economically comparable households. Any difference in effective rate can then be
attributed to the tax system changing between 1963, 1988, and 2013, rather than to the
households being different. It is what turns the comparison into a statement about tax
policy over time.
\end{solution}
\end{problem}

\section*{Part 2 --- More ``Backwards'' from Taxes}

Now each couple paid \textbf{\$3,000} in tax, and we work backwards to their income.
Same assumptions as before.

\begin{center}
\begin{tabular}{|l|c|c|c|}
\hline
\textbf{Description} & \textbf{1963} & \textbf{1988} & \textbf{2013} \\
\hline
Adjusted Gross Income (AGI) & & & \\
\hline
Standard deduction & \$1,200 & \$8,900 & \$20,000 \\
\hline
Taxable income & & & \\
\hline
Total tax & \$3,000 & \$3,000 & \$3,000 \\
\hline
Effective rate (\% of TOTAL income) & & & \\
\hline
\end{tabular}
\end{center}

\begin{problem}
Use the \textbf{1963} tax information to calculate how much taxable income a couple had
if they paid \$3,000 in taxes. Then, assuming the standard deduction shown, determine
their AGI.

\showworkyear

\begin{freeResponse}
\end{freeResponse}

\begin{solution}
Identify which row of the 1963 schedule could produce a \$3,000 tax --- each row covers a
known range of tax amounts, so find the row whose range contains \$3,000. Then reverse
that row's instruction: subtract the base amount, divide by the rate, and add the row's
lower boundary. Finally add back the \$1,200 deduction to recover the AGI.
\end{solution}
\end{problem}

\begin{problem}
Do the same with the \textbf{1988} tax information.

\showworkyear

\begin{freeResponse}
\end{freeResponse}

\begin{solution}
Same method against the 1988 schedule, adding back the \$8,900 deduction at the end.
\end{solution}
\end{problem}

\begin{problem}
Do the same with the \textbf{2013} tax information.

\showworkyear

\begin{freeResponse}
\end{freeResponse}

\begin{solution}
Same method against the 2013 schedule, adding back the \$20,000 deduction at the end.
\end{solution}
\end{problem}

\begin{problem}
Calculate the effective tax rate for each year of this second table, again as a
percentage of the \textbf{TOTAL} income (AGI).

\begin{freeResponse}
\end{freeResponse}

\begin{solution}
Divide \$3,000 by each year's recovered AGI. Since the numerator is now the same in all
three years, the effective rates are determined entirely by how much income it took to
generate a \$3,000 tax bill in that year.
\end{solution}
\end{problem}

\begin{problem}
In which year did paying \$3,000 in taxes ``cost the most''? Explain your choice.

\begin{freeResponse}
\end{freeResponse}

\begin{solution}
``Cost the most'' means the largest share of that household's income, so it is the year
with the highest effective rate --- equivalently, the year in which the \emph{smallest}
AGI was enough to generate a \$3,000 tax bill.
\end{solution}
\end{problem}

\begin{problem}
Why would we expect that year to be the answer even without having done any
calculations?

\begin{freeResponse}
\end{freeResponse}

\begin{solution}
Because \$3,000 is a fixed nominal amount, and inflation makes a dollar in the earliest
year worth far more than a dollar in the latest. In 1963 a \$3,000 tax bill is an
enormous sum relative to what people earned; by 2013 it is a modest one. So the earliest
year should require the least real income to produce that tax, and therefore show the
highest effective rate.

This is the mirror image of Part 1. There, we held the real income constant and let the
tax vary. Here, we hold the nominal tax constant and let the income vary --- so
inflation, rather than the structure of the schedules, drives most of the answer.
\end{solution}
\end{problem}

\end{document}
